\section{Conclusion}
\label{sec:conclusion}
On MathNet-Retrieve the two tiers reward two different things. The easy tier rewards the genre of LLM rewrites, whichever prompt wrote them; the hard tier rewards the recipe's pair structure, a deep rewrite set against a minimal-edit near-miss. Half to two thirds of the recipe arm's easy-tier lead over the verified arm is bought by any LLM-written training pairs, and back-translation, as deep as a non-LLM paraphrase reaches, attributes this to the authorship rather than the depth of the paraphrase; the rest needs the benchmark's own prompt and vanishes on real duplicates that no generator wrote. MELD also moves when trained on pairs built its way, SABER-Math only through a small gain robust to a matched budget, and only on MathNet-Retrieve can one edit to a training file be shown to raise the score and lower real retrieval at once. A benchmark should not make score and retention point in opposite directions.

Four practices follow for builders. (i)~Assume the recipe will get out or be rebuilt: MELD's score moved under a reconstruction of a prompt it never published. (ii)~Prefer test material that people wrote to material an LLM wrote. (iii)~Keep a subset no generator touched and score models on it too; a model that does much better on the LLM-written part is showing the gaming signal. (iv)~Publish a recipe-sensitivity number, how much a model gains by matching the recipe, as contamination tests publish how much of a score does not carry over to reworded or synthetic copies~\citep{dekoninck2024constat}. Regenerating the golds is a weaker lever: with the benchmark's own recipe it removes under a quarter of the gap, and drawing golds from several prompts under two fifths (Appendix~\ref{sec:forgetting}).
